\chapter{Zbiory analityczne i operacja Suslina}
\label{ch:topog-analityczne-suslin}

\noindent\textbf{Cel rozdziału.}
W rozdziale~\ref{ch:topog-polskie-borelowskie}
poznaliśmy przestrzenie polskie i zbiory borelowskie.
Teraz zbadamy zbiory powstające jako obrazy ciągłe
przestrzeni polskich i rzuty zbiorów domkniętych.
Wykażemy równoważność tych opisów, wprowadzimy
operację Suslina i skonstruujemy zbiór analityczny,
który nie jest borelowski.

Nie powtarzamy rozdziału~30 skryptu
\emph{Logika matematyczna i teoria mnogości},
poświęconego grom nieskończonym, aksjomatowi
determinacji, klasom projektywnym i redukcjom Wadge'a.
Podajemy natomiast pełne dowody topologicznych
faktów, które tam występują jedynie jako
wprowadzenie lub wyniki przywołane bez dowodu.
W całym rozdziale pracujemy w zwykłym
ujęciu teorii mnogości ZFC.

\section{Przestrzeń Baire'a jako przestrzeń parametryczna}
\label{sec:topog20-baire}

Przez \(\mathcal N=\mathbb N^{\mathbb N}\)
oznaczamy przestrzeń Baire'a z topologią
produktową, opisaną w
Przykładzie~\ref{ex:topog-polskie-cantor-baire}.
Indeksy współrzędnych zaczynamy od \(1\).
Zbiór wszystkich skończonych ciągów
liczb naturalnych oznaczamy
\(\mathbb N^{<\mathbb N}\).
Dla \(s=(s_1,\ldots,s_k)\) definiujemy
\[
 [s]=\{x\in\mathcal N:
             x_i=s_i\text{ dla }1\le i\le k\}.
\]
Przyjmujemy \([\varnothing]=\mathcal N\).
Zbiory \([s]\) są otwarte i domknięte,
a rodzina wszystkich \([s]\)
jest przeliczalną bazą topologii
\(\mathcal N\).
Dla \(x\in\mathcal N\) zapisujemy
\(x|k=(x_1,\ldots,x_k)\).

\begin{lemat}[Parametryzacja przestrzeni polskiej]
\label{lem:topog20-surjection}
Niech \(X\) będzie niepustą przestrzenią
polską. Istnieje ciągłe odwzorowanie
surjektywne \(f\colon\mathcal N\to X\).
\end{lemat}
\begin{proof}
Wybierzmy na \(X\) zgodną zupełną
metrykę \(d\) oraz przeliczalny
gęsty zbiór \(D\).
Zbudujemy rekurencyjnie niepuste
domknięte zbiory \(F_s\subseteq X\),
\(s\in\mathbb N^{<\mathbb N}\),
spełniające
\[
 F_{\varnothing}=X,\qquad
 F_{s^\frown k}\subseteq F_s,\qquad
 F_s=\bigcup_{k\in\mathbb N}F_{s^\frown k},
\]
oraz
\(\operatorname{diam}_d(F_s)\le2^{-|s|}\)
dla \(|s|\ge1\).

Załóżmy, że zdefiniowano
niepusty domknięty \(F_s\)
i niech \(l=|s|\).
Wszystkie niepuste zbiory postaci
\[
 F_s\cap\overline B_d(a,2^{-l-3}),
 \qquad a\in D,
\]
są domknięte, zawarte w \(F_s\)
i mają średnicę nie większą
niż \(2^{-l-2}\).
Pokrywają \(F_s\): dla \(x\in F_s\)
z gęstości \(D\) można wybrać
\(a\in D\) z \(d(x,a)<2^{-l-3}\).
Tę niepustą przeliczalną rodzinę
enumerujemy jako
\((F_{s^\frown k})_{k\in\mathbb N}\);
jeśli jest skończona,
powtarzamy jej elementy.
To kończy konstrukcję indukcyjną.

Dla \(\alpha\in\mathcal N\)
zbiory \(F_{\alpha|k}\)
tworzą malejącą rodzinę
niepustych domkniętych
podzbiorów zupełnej przestrzeni
metrycznej, a ich średnice
dążą do zera.
Wybierzmy \(x_k\in F_{\alpha|k}\).
Dla \(p,q\ge k\) punkty
\(x_p,x_q\) należą do
\(F_{\alpha|k}\), zatem
\(d(x_p,x_q)\le2^{-k}\).
Ciąg \((x_k)\) jest więc
ciągiem Cauchy'ego i ma
granicę \(x\in X\).
Domkniętość każdego
\(F_{\alpha|k}\) zapewnia
\(x\in F_{\alpha|k}\).
Średnice pokazują zarazem
jednoznaczność takiego \(x\).
Definiujemy \(f(\alpha)=x\).

Jeżeli \(\alpha,\beta\)
mają ten sam początkowy
fragment długości \(k\), to
\(f(\alpha)\) i \(f(\beta)\)
leżą w \(F_{\alpha|k}\),
więc ich odległość
nie przekracza \(2^{-k}\).
Stąd \(f\) jest ciągłe.

Dla dowolnego \(x\in X\)
wybieramy rekurencyjnie
\(k_1,k_2,\ldots\), aby
\(x\in F_{(k_1,\ldots,k_n)}\)
dla wszystkich \(n\).
Wybór jest możliwy, bo
każdy \(F_s\) jest sumą
swoich następców.
Dla \(\alpha=(k_n)\)
otrzymujemy \(f(\alpha)=x\).
Zatem \(f\) jest surjektywne.
\end{proof}

\begin{uwaga}[Co daje parametryzacja]
\label{rem:topog20-param}
W powyższym lemacie \(X\)
nie musi być zwartą ani
lokalnie zwartą przestrzenią.
Nie twierdzimy również,
że dowolny ciągły obraz
przestrzeni polskiej jest polski.
Właśnie to rozróżnienie
prowadzi do zbiorów analitycznych.
\end{uwaga}

\section{Zbiory analityczne i rzuty zbiorów domkniętych}
\label{sec:topog20-analityczne}

\begin{definicja}[Zbiór analityczny]
\label{def:topog20-analityczny}
Niech \(X\) będzie przestrzenią polską.
Podzbiór \(A\subseteq X\) nazywamy
\emph{analitycznym}, jeśli
\(A=\varnothing\) albo istnieje
ciągłe odwzorowanie
\(f\colon\mathcal N\to X\)
takie, że \(A=f(\mathcal N)\).
Dopełnienie zbioru analitycznego
nazywamy \emph{koanalitycznym}.
\end{definicja}

\begin{stwierdzenie}[Inna przestrzeń parametryczna]
\label{prop:topog20-param}
Dla przestrzeni polskiej \(X\)
i niepustego \(A\subseteq X\)
następujące warunki są równoważne:
\begin{enumerate}[label=\textup{(\roman*)}]
\item \(A\) jest analityczny;
\item istnieje niepusta przestrzeń
polska \(Y\) i ciągłe odwzorowanie
\(g\colon Y\to X\) takie, że \(g(Y)=A\).
\end{enumerate}
\end{stwierdzenie}
\begin{proof}
Implikację \textup{(i)} \(\Rightarrow\)
\textup{(ii)} uzyskujemy,
biorąc \(Y=\mathcal N\).
Odwrotnie,
z Lematu~\ref{lem:topog20-surjection}
istnieje ciągła surjekcja
\(h\colon\mathcal N\to Y\).
Wówczas \(g\circ h\)
jest ciągłe i ma obraz \(A\).
\end{proof}

\begin{twierdzenie}[Charakterystyka przez rzut]
\label{thm:topog20-rzut}
Niech \(X\) będzie przestrzenią polską,
a \(A\subseteq X\).
Wtedy \(A\) jest analityczny
dokładnie wtedy, gdy istnieje
zbiór domknięty
\(F\subseteq X\times\mathcal N\)
taki, że
\[
 A=\pi_X(F)
   =\{x\in X:
       \text{ istnieje }u\in\mathcal N
       \text{ z }(x,u)\in F\}.
\]
\end{twierdzenie}
\begin{proof}
Jeżeli \(A=\varnothing\),
wystarczy przyjąć \(F=\varnothing\).
W przeciwnym przypadku,
jeśli \(A=f(\mathcal N)\)
dla ciągłego \(f\),
rozważamy jego graf
\[
 F=\{(f(u),u):u\in\mathcal N\}.
\]
Graf jest domknięty:
jeśli \((x_j,u_j)\to(x,u)\)
i \(x_j=f(u_j)\),
to z ciągłości \(f\)
wynika \(x=f(u)\).
Jego rzutem na \(X\)
jest dokładnie \(A\).

W drugą stronę załóżmy,
że \(A=\pi_X(F)\)
dla domkniętego \(F\).
Produkt \(X\times\mathcal N\)
jest przestrzenią polską na mocy
Twierdzenia~\ref{thm:topog-polskie-iloczyn}
(uzupełniając iloczyn dwoma czynnikami
o jednoelementowe przestrzenie polskie).
Jego domknięta podprzestrzeń
\(F\) jest polska.
Jeśli \(F=\varnothing\), to
\(A=\varnothing\).
W przeciwnym przypadku
Lemat~\ref{lem:topog20-surjection}
daje ciągłą surjekcję
\(h\colon\mathcal N\to F\).
Złożenie \( \pi_X|_F\circ h \)
jest ciągłe i ma obraz \(A\).
\end{proof}

\begin{przyklad}[Dlaczego rzut nie musi być domknięty]
\label{ex:topog20-rzut}
W \(\mathbb R^2\) zbiór
\[
 F=\{(x,y)\in\mathbb R^2:xy=1\}
\]
jest domknięty, bo jest
przeciwobrazem singletonu
\(\{1\}\) przez ciągłe
mnożenie. Tymczasem
\(\pi_1(F)=\mathbb R\setminus\{0\}\)
nie jest domknięty.
Przykład nie daje zbioru
nieborelowskiego: pokazuje tylko,
dlaczego własności domkniętości
nie wolno przenosić przez
dowolne rzutowanie.
\end{przyklad}

\section{Działania na zbiorach analitycznych}
\label{sec:topog20-operacje}

\begin{twierdzenie}[Przeliczalne sumy i przecięcia]
\label{thm:topog20-closure}
Niech \(X\) będzie przestrzenią polską
i niech \((A_n)_{n\ge1}\)
będzie ciągiem jej podzbiorów
analitycznych.
Wówczas zbiory
\(\bigcup_{n\ge1}A_n\) i
\(\bigcap_{n\ge1}A_n\)
są analityczne.
\end{twierdzenie}
\begin{proof}
Z Twierdzenia~\ref{thm:topog20-rzut}
dla każdego \(n\) możemy wybrać
domknięty \(F_n\subseteq
X\times\mathcal N\) taki,
że \(\pi_X(F_n)=A_n\).

Najpierw rozważmy przestrzeń
\(\mathbb N\times\mathcal N\),
gdzie \(\mathbb N\) jest dyskretna.
Jest ona homeomorficzna
z \(\mathcal N\), ponieważ
\((n,u_1,u_2,\ldots)\)
można utożsamić z
\((n,u_1,u_2,\ldots)\in\mathcal N\).
Zbiór
\[
 F=\{(x,(n,u)):
               (x,u)\in F_n\}
 \subseteq X\times(\mathbb N\times\mathcal N)
\]
jest domknięty: jego przekrój
nad każdym otwarto-domkniętym
\(\{n\}\) jest domknięty,
a dopełnienie \(F\) jest
sumą otwartych dopełnień
tych przekrojów.
Jego rzut na \(X\) jest
równy \(\bigcup_n A_n\).
Twierdzenie~\ref{thm:topog20-rzut}
daje analityczność sumy.

Dla przecięcia użyjmy
\(\mathcal N^{\mathbb N}\),
która jest homeomorficzna
z \(\mathcal N\): wystarczy
ustalić bijekcję
\(\mathbb N^2\to\mathbb N\)
i ponumerować wszystkie
współrzędne obu ciągów indeksów.
Definiujemy
\[
 H=\left\{(x,(u_n)_{n\ge1}):
            (x,u_n)\in F_n
            \text{ dla każdego }n\right\}.
\]
Jest to przecięcie
przeciwobrazów domkniętych \(F_n\)
przez ciągłe odwzorowania
\((x,(u_j))\mapsto(x,u_n)\);
jest zatem domknięty.
Jego rzutem jest
\(\bigcap_n A_n\):
w jednym kierunku wynika to
bezpośrednio z definicji,
w drugim dla \(x\in\bigcap_n A_n\)
wybieramy świadków \(u_n\)
z \((x,u_n)\in F_n\).
Ponownie stosujemy
Twierdzenie~\ref{thm:topog20-rzut}.
\end{proof}

\begin{stwierdzenie}[Obrazy i przeciwobrazy]
\label{prop:topog20-maps}
Niech \(X,Y\) będą przestrzeniami
polskimi i \(f\colon X\to Y\)
odwzorowaniem ciągłym.
Wówczas:
\begin{enumerate}[label=\textup{(\roman*)}]
\item jeżeli \(A\subseteq X\)
jest analityczny, to \(f(A)\)
jest analityczny w \(Y\);
\item jeżeli \(B\subseteq Y\)
jest analityczny, to
\(f^{-1}(B)\) jest analityczny w \(X\).
\end{enumerate}
\end{stwierdzenie}
\begin{proof}
Jeśli \(A=\varnothing\),
pierwsza teza jest oczywista.
Dla \(A=g(\mathcal N)\)
mamy \(f(A)=(f\circ g)(\mathcal N)\).

Dla drugiej tezy przedstawmy
\(B=\pi_Y(F)\), gdzie
\(F\subseteq Y\times\mathcal N\)
jest domknięty.
Odwzorowanie
\[
 \Phi\colon X\times\mathcal N
     \longrightarrow Y\times\mathcal N,
 \qquad \Phi(x,u)=(f(x),u)
\]
jest ciągłe.
Zbiór \(\Phi^{-1}(F)\)
jest więc domknięty,
a jego rzut na \(X\)
jest równy \(f^{-1}(B)\).
\end{proof}

\begin{wniosek}[Zbiory borelowskie są analityczne]
\label{cor:topog20-borel}
W każdej przestrzeni polskiej \(X\)
każdy zbiór borelowski \(B\subseteq X\)
jest zarazem analityczny i koanalityczny.
\end{wniosek}
\begin{proof}
Niech \(\mathcal D\)
oznacza rodzinę wszystkich
zbiorów \(B\subseteq X\),
dla których zarówno \(B\),
jak i \(X\setminus B\)
są analityczne.
Na mocy definicji
\(\mathcal D\) jest zamknięta
na dopełnienia.
Twierdzenie~\ref{thm:topog20-closure}
oraz prawa de Morgana
pokazują, że jest zamknięta
na przeliczalne sumy;
jest więc \(\sigma\)-algebrą.

Każdy zbiór otwarty \(U\)
należy do \(\mathcal D\):
zarówno \(U\), jak i
jego dopełnienie są borelowskie
na poziomie pierwszym,
ale analityczność uzasadniamy
bez odwołania do tezy.
W przestrzeni metrycznej
\(U\) jest typu \(F_\sigma\)
na mocy Lematu~
\ref{lem:topog-polskie-otwarte-fsigma}.
Każdy zbiór domknięty,
jako podprzestrzeń polska,
jest analityczny z
Lematu~\ref{lem:topog20-surjection}
(lub jest pusty).
Stąd \(U\), jako przeliczalna
suma takich zbiorów, jest
analityczny. Jego dopełnienie
jest domknięte, więc również
analityczne.
Ponieważ \(\mathcal D\)
zawiera wszystkie otwarte
i jest \(\sigma\)-algebrą,
zawiera całą
\(\mathcal B(X)\).
\end{proof}


\begin{stwierdzenie}[Zgodność z definicją przez rzut borelowski]
\label{prop:topog20-borel-projection}
Niech \(n\ge1\), a \(A\subseteq\mathbb R^n\).
Zbiór \(A\) jest analityczny
w sensie Definicji~\ref{def:topog20-analityczny}
dokładnie wtedy, gdy
istnieje zbiór borelowski
\(B\subseteq\mathbb R^{n+1}\)
taki, że \(A=\pi_{\mathbb R^n}(B)\).
Jest to definicja stosowana
w rozdziale~30 skryptu
\emph{Logika matematyczna i teoria mnogości}.
\end{stwierdzenie}
\begin{proof}
Jeżeli \(B\) jest borelowski
w \(\mathbb R^{n+1}\),
to jest analityczny na mocy
Wniosku~\ref{cor:topog20-borel}.
Rzutowanie
\(\mathbb R^{n+1}\to\mathbb R^n\)
jest ciągłe, zatem
Stwierdzenie~\ref{prop:topog20-maps}
pokazuje, że jego obraz \(A\)
jest analityczny.

Dla implikacji odwrotnej
skonstruujemy homeomorficzne
zanurzenie przestrzeni Baire'a
w borelowski podzbiór prostej.
Niech \(\mathbb N_0=\{0,1,2,\ldots\}\)
i ustalmy bijekcję
\(b\colon\mathbb N\to\mathbb N_0\).
Dla \(\alpha=(\alpha_k)\in\mathcal N\)
zapiszmy nieskończony ciąg binarny
\[
 e(\alpha)=
 0^{b(\alpha_1)}1\,
 0^{b(\alpha_2)}1\,
 0^{b(\alpha_3)}1\cdots,
\]
gdzie \(0^l\) oznacza blok
złożony z \(l\) zer,
a \(0^0\) jest blokiem pustym.
Obrazem \(e\) jest zbiór
\[
 C_\infty=\{u\in2^{\mathbb N}:
             u_k=1\text{ dla
             nieskończenie wielu }k\}.
\]
Jest on typu \(G_\delta\)
w przestrzeni Cantora, ponieważ
\[
 C_\infty=
 \bigcap_{m\ge1}\bigcup_{k\ge m}
                  \{u:u_k=1\}.
\]
Odwzorowanie \(e\) jest bijekcją
na \(C_\infty\).
Jest ciągłe, bo początkowy
fragment jego wartości
zależy tylko od skończenie
wielu współrzędnych \(\alpha\).
Jego odwrotność jest ciągła:
pierwszych \(l\) wyrazów
\(\alpha\) odczytujemy z długości
bloków zer przed pierwszymi
\(l\) jedynkami, a ich wystąpienie
jest rozstrzygnięte przez
skończony fragment ciągu binarnego.

Odwzorowanie
\[
 j\colon2^{\mathbb N}\longrightarrow
           \mathbb R,\qquad
 j(u)=\sum_{k\ge1}\frac{2u_k}{3^k}
\]
jest ciągłe i różnowartościowe:
jeśli dwa ciągi różnią się
po raz pierwszy na pozycji \(k\),
różnica wkładu tego wyrazu
ma moduł \(2/3^k\),
a cały ogon ma moduł
co najwyżej \(1/3^k\).
Ponieważ \(2^{\mathbb N}\)
jest zwarta, \(j\) jest
homeomorfizmem na swój obraz,
który jest domknięty w
\(\mathbb R\).
Zatem \(E=j(C_\infty)\)
jest borelowski w \(\mathbb R\),
a \(i=j\circ e\)
jest homeomorfizmem
\(\mathcal N\) na \(E\).

Niech teraz \(A\) będzie analityczny.
Z Twierdzenia~\ref{thm:topog20-rzut}
mamy domknięty
\(F\subseteq\mathbb R^n\times\mathcal N\)
z \(\pi_{\mathbb R^n}(F)=A\).
Obraz \(F\) przez
\(\operatorname{id}_{\mathbb R^n}\times i\)
jest domknięty w przestrzeni
względnej \(\mathbb R^n\times E\).
Ponieważ \(E\) jest borelowski,
istnieje domknięty
\(H\subseteq\mathbb R^{n+1}\),
dla którego ten obraz
jest równy
\(H\cap(\mathbb R^n\times E)\).
Jest on zatem borelowski
w \(\mathbb R^{n+1}\)
i ma rzut równy \(A\).
\end{proof}

\section{Operacja Suslina}
\label{sec:topog20-suslin}

\begin{definicja}[Operacja Suslina]
\label{def:topog20-suslin}
Niech \(X\) będzie zbiorem,
a \((F_s)_{s\in\mathbb N^{<\mathbb N}}\)
rodziną jego podzbiorów.
\emph{Operacją Suslina}
na tej rodzinie nazywamy zbiór
\[
 \mathcal A(F_s)=
 \bigcup_{\alpha\in\mathcal N}
     \bigcap_{n\ge1}F_{\alpha|n}.
\]
W definicji nie wymagamy
monotoniczności rodziny \(F_s\).
\end{definicja}

\begin{twierdzenie}[Charakterystyka Suslina]
\label{thm:topog20-suslin}
Niech \(X\) będzie przestrzenią polską,
a \(A\subseteq X\).
Następujące warunki są równoważne:
\begin{enumerate}[label=\textup{(\roman*)}]
\item \(A\) jest zbiorem analitycznym;
\item istnieje rodzina domkniętych
zbiorów \(F_s\subseteq X\),
indeksowana wszystkimi
\(s\in\mathbb N^{<\mathbb N}\),
taka, że \(A=\mathcal A(F_s)\).
\end{enumerate}
\end{twierdzenie}
\begin{proof}
Przypuśćmy najpierw, że
\(A=\mathcal A(F_s)\)
dla domkniętych \(F_s\).
Określmy
\[
 H=\{(x,\alpha)\in X\times\mathcal N:
       x\in F_{\alpha|n}
       \text{ dla każdego }n\ge1\}.
\]
Dla ustalonego \(n\) zbiór
\[
 H_n=\bigcup_{s\in\mathbb N^n}
                         (F_s\times[s])
\]
jest domknięty.
Istotnie, cylindry \([s]\)
o długości \(n\) tworzą
otwarto-domknięty podział
\(\mathcal N\), a jego dopełnienie
jest sumą otwartych zbiorów
\((X\setminus F_s)\times[s]\).
Mamy \(H=\bigcap_n H_n\),
zatem \(H\) jest domknięty.
Jego rzut na \(X\)
jest równy \(\mathcal A(F_s)\),
więc z
Twierdzenia~\ref{thm:topog20-rzut}
wynika analityczność \(A\).

Odwrotnie, niech
\(A=\pi_X(H)\) dla domkniętego
\(H\subseteq X\times\mathcal N\).
Wybierzmy zgodną metrykę \(d\)
na \(X\) i dla skończonego \(s\)
połóżmy
\[
 F_s=
 \overline{\pi_X\bigl(
     H\cap(X\times[s])\bigr)}^{\,X}.
\]
Są to zbiory domknięte.
Jeśli \((x,\alpha)\in H\),
to \(x\in F_{\alpha|n}\)
dla każdego \(n\),
więc \(A\subseteq\mathcal A(F_s)\).

Niech teraz
\(x\in\mathcal A(F_s)\).
Istnieje \(\alpha\in\mathcal N\)
takie, że
\(x\in F_{\alpha|n}\)
dla każdego \(n\).
Z definicji domknięcia
możemy dla każdego \(n\)
wybrać \((x_n,u_n)\in H\)
tak, aby
\[
 d(x_n,x)<1/n,\qquad
 u_n\in[\alpha|n].
\]
Zatem \(x_n\to x\),
a \(u_n\to\alpha\)
w topologii produktowej.
Ponieważ \(H\) jest domknięty
w produkcie przestrzeni metryzowalnych,
otrzymujemy \((x,\alpha)\in H\).
W szczególności
\(x\in\pi_X(H)=A\).
Dowodzi to drugiego zawierania.
\end{proof}

\begin{uwaga}[Granica ustalonego zakresu]
\label{rem:topog20-no-meas}
Powyższe twierdzenie daje
topologiczny opis zbiorów analitycznych.
W szczególności \emph{nie} stanowi
jeszcze dowodu ich uniwersalnej
mierzalności ani twierdzenia
o rozdzielaniu zbiorów analitycznych.
W skrypcie \emph{Logika matematyczna
i teoria mnogości}, rozdział~30,
twierdzenie o mierzalności
zbiorów analitycznych jest
jawnie przywołane jako zewnętrzny wynik.
Nie przedstawiamy go tutaj
jako udowodnionego.
\end{uwaga}

\section{Zbiór analityczny nieborelowski}
\label{sec:topog20-nieborelowski}

\begin{lemat}[Uniwersalna rodzina rzutów domkniętych]
\label{lem:topog20-universal}
Istnieje analityczny
\(U\subseteq\mathcal N\times\mathcal N\)
o następującej własności:
dla każdego analitycznego
\(A\subseteq\mathcal N\)
istnieje \(c\in\mathcal N\)
takie, że
\[
 A=U_c=\{x\in\mathcal N:(c,x)\in U\}.
\]
\end{lemat}
\begin{proof}
Zbiór par \((s,t)\) skończonych
ciągów naturalnych tej samej
długości jest przeliczalny.
Ustalmy jego numerację
\(e(s,t)\in\mathbb N\), będącą bijekcją.
Dla \(c\in\mathcal N\) kodujemy
rodzinę par przez warunek
\(c_{e(s,t)}=1\).
Wartość \(1\) oznacza obecność pary,
a wartość \(2\) jej brak.
Obie wartości należą do \(\mathbb N\),
niezależnie od tego, czy \(0\)
zaliczamy do liczb naturalnych.

Zdefiniujmy
\[
 H=\{(c,x,y)\in\mathcal N^3:
       c_{e(x|n,y|n)}=1
       \text{ dla każdego }n\ge1\}.
\]
Jest to zbiór domknięty.
Dla ustalonego \(n\)
warunek na pierwszych
\(n\) współrzędnych \(x,y\)
i jedną współrzędną \(c\)
wyznacza zbiór otwarty
i domknięty; przecięcie
tych zbiorów po \(n\)
jest domknięte.
Niech
\[
 U=\{(c,x):\text{ istnieje }y
                   \text{ z }(c,x,y)\in H\}.
\]
Z Twierdzenia~\ref{thm:topog20-rzut}
\(U\) jest analityczny,
po utożsamieniu
\(\mathcal N^2\cong\mathcal N\).

Weźmy teraz dowolny analityczny
\(A\subseteq\mathcal N\).
Istnieje domknięty
\(F\subseteq\mathcal N^2\)
z \(\pi_1(F)=A\).
Rozważmy rodzinę
\[
 T_F=\{(s,t):
    |s|=|t|,\quad
    F\cap([s]\times[t])
           \ne\varnothing\}.
\]
Dobierzmy kod \(c\) tak,
że \(c_{e(s,t)}=1\) dla par
\((s,t)\in T_F\), natomiast
\(c_{e(s,t)}=2\) dla
wszystkich pozostałych par.

Dla \((x,y)\in F\)
wszystkie prefiksy
\((x|n,y|n)\) należą do \(T_F\).
Odwrotnie, jeśli każdy
taki prefiks należy do \(T_F\),
dla każdego \(n\) wybieramy
\((x_n,y_n)\in
F\cap([x|n]\times[y|n])\).
Wtedy \((x_n,y_n)\to(x,y)\),
a domkniętość \(F\) daje
\((x,y)\in F\).
Zatem
\[
 (c,x,y)\in H
   \ \Longleftrightarrow\ (x,y)\in F,
\]
co po rzutowaniu na \(x\)
daje \(U_c=A\).
\end{proof}

\begin{twierdzenie}[Istnienie nieborelowskiego zbioru analitycznego]
\label{thm:topog20-nonborel}
W przestrzeni Baire'a
\(\mathcal N=\mathbb N^{\mathbb N}\)
istnieje zbiór analityczny,
który nie jest borelowski.
\end{twierdzenie}
\begin{proof}
Niech \(U\) będzie zbiorem
z Lematu~\ref{lem:topog20-universal}.
Określmy
\[
 D=\{x\in\mathcal N:
                 (x,x)\notin U\}.
\]
Przypuśćmy, że \(D\)
jest analityczny.
Z uniwersalności \(U\)
istnieje \(c\in\mathcal N\)
takie, że \(D=U_c\).
Wtedy
\[
 c\in D
 \ \Longleftrightarrow\ (c,c)\notin U
 \ \Longleftrightarrow\ c\notin U_c=D,
\]
co jest sprzecznością.
Zatem \(D\) nie jest analityczny.

Jego dopełnienie
\[
 E=\{x\in\mathcal N:(x,x)\in U\}
\]
jest analityczne:
to przeciwobraz analitycznego
\(U\) przez ciągłe
odwzorowanie diagonalne
\(x\mapsto(x,x)\), zgodnie ze
Stwierdzeniem~\ref{prop:topog20-maps}.
Gdyby \(E\) był borelowski,
jego dopełnienie \(D\)
też byłoby borelowskie,
a więc analityczne na mocy
Wniosku~\ref{cor:topog20-borel}.
Otrzymujemy sprzeczność.
Wobec tego \(E\) jest
analityczny i nieborelowski.
\end{proof}

\begin{wniosek}[Rzut może zwiększać złożoność]
\label{cor:topog20-projection}
Istnieje zbiór domknięty
\(F\subseteq\mathcal N^2\),
którego rzut na pierwszą
współrzędną nie jest borelowski.
\end{wniosek}
\begin{proof}
Zbiór \(E\) otrzymany w
Twierdzeniu~\ref{thm:topog20-nonborel}
jest analityczny i nieborelowski.
Z Twierdzenia~\ref{thm:topog20-rzut}
jest rzutem zbioru domkniętego
\(F\subseteq\mathcal N^2\).
\end{proof}

\section{Podsumowanie i granice materiału}
\label{sec:topog20-summary}

Niech \(X\) będzie przestrzenią polską.
Podzbiór \(X\) jest analityczny
dokładnie wtedy, gdy jest pusty
albo jest ciągłym obrazem
przestrzeni Baire'a.
Równoważnie, jest rzutem
domkniętego podzbioru
\(X\times\mathcal N\) lub wynikiem
operacji Suslina na rodzinie
zbiorów domkniętych.
Są zamknięte na przeliczalne
sumy, przeliczalne przecięcia,
ciągłe obrazy i przeciwobrazy,
a każdy zbiór borelowski
jest analityczny i koanalityczny.
W przestrzeni Baire'a
klasa zbiorów analitycznych
jest ściśle większa od klasy
zbiorów borelowskich, zgodnie z
Twierdzeniem~\ref{thm:topog20-nonborel}.
Nie zachodzi to dla każdej
przestrzeni polskiej:
w przeliczalnej przestrzeni
dyskretnej każdy podzbiór
jest otwarty, a zatem borelowski.

Nie powtarzaliśmy teorii
determinacji i redukcji Wadge'a
z rozdziału~30 skryptu
\emph{Logika matematyczna i teoria mnogości}.
Wyniki dotyczące
uniwersalnej mierzalności,
własności zbioru doskonałego
i rozdzielania zbiorów analitycznych
pozostają osobnymi,
głębszymi zagadnieniami.
